\section{Built for the Rich Parts of the World}
\label{sec:10.5}
AI development can narrow global inequality or widen it, and which one happens is not automatic. A model trained overwhelmingly on data from wealthy, English-speaking, urban contexts underperforms everywhere else, and a government or company with no reason to fix that gap usually does not.

\begin{esbox}
\boxtitle{The scale of the connectivity gap, 2025}
The International Telecommunication Union's 2025 assessment put global internet use at roughly six billion people, about three-quarters of the world's population, with 2.2 billion people offline, concentrated overwhelmingly in low- and middle-income countries. The gap was starker in quality than in bare access: 94 percent of people in high-income countries used the internet against 23 percent in low-income ones, and 5G coverage reached 84 percent of high-income populations against 4 percent of low-income ones \autocite{itu2025facts}.

\end{esbox}
The constraints behind those numbers are structural. Reliable data, computing infrastructure, and the specialized personnel needed to build and maintain AI systems are unevenly distributed, and a country short on all three cannot simply adopt a system built elsewhere and expect it to transfer cleanly. Weak regulatory capacity compounds the problem: the states least equipped to build AI are correspondingly the states least equipped to govern whatever gets deployed within their borders, regardless of who built it. The export controls discussed above cut both ways here, since a control written to stop a surveillance system from reaching an authoritarian buyer can as easily gatekeep the same technology away from a legitimate one.

The remedies proposed at this level have been the same handful for years, every one of them is correct, none is contested, and they have been repeated in successive declarations the whole time — most recently the Global Digital Compact, adopted by all 193 UN member states in 2024, which created two follow-on bodies and no obligation \autocite{un2024globaldigitalcompact}. The box is what happened anyway. A government or company treating AI primarily as strategic advantage has no reason to fund any of them, and a declaration is not a reason.

\runin{What the geopolitics leaves to design}
A system that resists authoritarian capture still assumes a government somewhere is deciding whether to let it be built, sold, or deployed as designed, and those governments are unevenly equipped. None of that is a reason to abandon the design-level arguments — a system built to resist misuse resists it whichever government asks — but it is a reason not to expect law, international or domestic, to do work only a design can do.
